\section{Same-object Artin character factors}
\label{sec:factorization}

\subsection{Regular and isotypic transfers}

Let \(L_a=\left(\begin{smallmatrix}0&1\\1&0\end{smallmatrix}\right)\) denote
left translation by the nonidentity element of \(C_2\).  The frozen regular
transfer is
\[
B_{\mathrm{reg},P}(x)=
\sum_{\varnothing\ne S\subseteq P}\eps(S)x_SL_a^{|S|},
\qquad
D_{\mathrm{reg},P}(x)=\det(I-B_{\mathrm{reg},P}(x)).
\tag{4.1}\label{4.1}
\]
For the characters \(\chi_+(a)=1\) and \(\chi_-(a)=-1\), set
\[
B_{\chi,P}(x)=\sum_S\eps(S)x_S\chi(a)^{|S|},
\qquad D_{\chi,P}(x)=1-B_{\chi,P}(x).
\tag{4.2}\label{4.2}
\]

\begin{theorem}[Finite-group Artin decomposition]
\label{thm:artin}
Let \(G\) be finite and let \(\alpha:\E_P\to G\) be any one-letter cocycle.
For
\[
\Breg=\sum_S\eps(S)x_SL_{\alpha(S)}
\]
and
\[
B_\rho=\sum_S\eps(S)x_S\rho(\alpha(S)),
\qquad D_\rho=\det(I-B_\rho),
\]
one has
\[
\det(I-\Breg)=\prod_{\rho\in\whG}D_\rho^{d_\rho},
\qquad d_\rho=\dim\rho.
\tag{4.3}\label{4.3}
\]
The trivial block is independent of the cocycle and equals
\[
D_{\one}(x)=\prod_{p\in P}(1-x_p).
\tag{4.4}\label{4.4}
\]
\end{theorem}

\begin{proof}
Left convolution commutes with right deck translations.  The regular
representation decomposes into \(d_\rho\) copies of each irreducible
\(\rho\), so determinants multiply blockwise.  Equation \eqref{4.4} is finite
inclusion--exclusion.  Full details appear in \cref{app:proof-artin}.
\end{proof}

Theorem~\ref{thm:artin} is classical machinery, not the novelty claim.  Its
role is logical: every factor below belongs to the same regular transfer.

\subsection{Degree-power factorization}

\begin{theorem}[Degree-power atom locality]
\label{thm:degree-power}
Fix \(a\in G\), set \(\alpha(S)=a^{|S|}\), and write \(A=\rho(a)\).  Then
\[
I-B_\rho(x)=\prod_{p\in P}(I-x_pA),
\qquad
D_\rho(x)=\prod_{p\in P}\det(I-x_pA).
\tag{4.5}\label{4.5}
\]
\end{theorem}

\begin{proof}
Expanding the commuting matrix product gives
\[
\prod_p(I-x_pA)=
I+\sum_{k\ge1}(-1)^ke_k(x)A^k
=I-B_\rho(x).
\]
Taking determinants proves the second identity.
\end{proof}

For the primary \(C_2\) extension, \cref{thm:degree-power} gives
\[
D_+(x)=\prod_p(1-x_p),\qquad
D_-(x)=\prod_p(1+x_p),\qquad
\Dreg(x)=\prod_p(1-x_p^2).
\tag{4.6}\label{4.6}
\]
At \(x_p=p^{-s}\), \(\Ree s>1\),
\[
D_+(s)=\frac1{\zeta(s)},\qquad
D_-(s)=\frac{\zeta(s)}{\zeta(2s)},\qquad
\Dreg(s)=\frac1{\zeta(2s)}.
\tag{4.7}\label{4.7}
\]

\begin{figure}[t]
\centering
\resizebox{\textwidth}{!}{%
\begin{tikzpicture}[
  node distance=8mm and 10mm,
  box/.style={draw,rounded corners=2pt,align=center,minimum height=12mm,
              inner sep=5pt,line width=.7pt},
  source/.style={box,draw=fiberblue,fill=fiberblue!6},
  block/.style={box,draw=blockgreen,fill=blockgreen!7},
  whole/.style={box,draw=routegray,fill=routegray!6,line width=1pt},
  stop/.style={box,draw=stopamber,fill=stopamber!8,dashed},
  arr/.style={-{Latex[length=2.2mm]},line width=.7pt}
]
\node[source] (source) {signed subset shift\\
  \(w(S)=(-1)^{|S|+1}x_S\)\\
  \(\alpha(S)=|S|\bmod2\)};
\node[source,right=of source] (regular) {\(C_2\) regular transfer\\
  \(\displaystyle \Breg=\sum_S\eps(S)x_SL_a^{|S|}\)\\
  deck action commutes};
\node[block,above right=2mm and 12mm of regular] (plus) {trivial isotype\\
  \(D_+=\prod_p(1-x_p)\)};
\node[block,below right=2mm and 12mm of regular] (minus) {sign isotype\\
  \(D_-=\prod_p(1+x_p)\)};
\node[whole,right=15mm of regular] (whole) {whole extension\\
  \(\Dreg=D_+D_-\)\\
  \(\displaystyle =\prod_p(1-x_p^2)\)};
\draw[arr] (source) -- (regular);
\draw[arr] (regular.east) -- ++(5mm,0) |- (plus.west);
\draw[arr] (regular.east) -- ++(5mm,0) |- (minus.west);
\draw[arr] (plus.east) -- ++(5mm,0) |- (whole.west);
\draw[arr] (minus.east) -- ++(5mm,0) |- (whole.west);

\node[stop,below=14mm of source,xshift=13mm] (primitive) {primitive stop\\
  singleton clock \(\times m\); mixed lifts remain};
\node[stop,right=12mm of primitive] (controls) {selectivity stop\\
  arbitrary inventories reproduce the identity};
\node[whole,right=12mm of controls] (route) {\(\RouteA\) rejected\\
  \(\RouteB\) locked};
\draw[arr,dashed,draw=stopamber] (regular.south) |- (primitive.north);
\draw[arr,dashed,draw=stopamber] (whole.south) |- (controls.north);
\draw[arr] (primitive) -- (controls);
\draw[arr] (controls) -- (route);
\end{tikzpicture}}
\caption{Same-object Artin factorization and its scope.  Both character
determinants are isotypic blocks of one regular transfer; their product alone
is the whole-extension determinant.  Atom-local factors coexist with mixed
lifted primitives and exact replication by nonprime controls.}
\label{fig:factorization-scope}
\end{figure}


\paragraph{Whole versus block.}
Only \(\Dreg\) in \eqref{4.6} is the determinant of the whole regular
extension.  The quantities \(D_+\) and \(D_-\) are determinants of invariant
isotypic blocks of that same operator.  Calling either block the
whole-extension determinant would erase the Artin decomposition rather than
use it.

\paragraph{Two-atom certificate.}
For variables \(x,y\), the singleton symbols have weights \(x,y\) and fiber
label \(a\), while \(\{p,q\}\) has weight \(-xy\) and identity fiber label.
Thus
\[
\Breg=(x+y)L_a-xyI,\qquad
\det(I-\Breg)=(1-x^2)(1-y^2),
\tag{4.8}\label{4.8}
\]
and diagonalizing \(L_a\) yields
\[
D_+=(1-x)(1-y),\qquad D_-=(1+x)(1+y).
\tag{4.9}\label{4.9}
\]
This is the smallest direct certificate that motion, character factors, and
the whole cover use one matrix and one normalization.

\paragraph{General cyclic fiber.}
For \(C_m=\langle a\rangle\), \(\omega=e^{2\pi i/m}\), and
\(\chi_j(a)=\omega^j\),
\[
D_j(x)=\prod_p(1-\omega^jx_p),\qquad
\Dreg(x)=\prod_{j=0}^{m-1}D_j(x)=\prod_p(1-x_p^m).
\tag{4.10}\label{4.10}
\]
The local factors in \eqref{4.5} and \eqref{4.10} are indexed by single
atoms.  This does not say that their coefficient expansions lack mixed
monomials or that the symbolic extension lacks mixed primitive cycles.
